% Algorithmic Non-commutative Class Field Theory, IV
% The limit module of abelian towers and the characteristic identity
% Build: pdflatex anccft4 (x3)
\documentclass[11pt]{amsart}
\usepackage{amsmath,amssymb,amsthm,mathtools}
\usepackage[margin=1.0in]{geometry}
\usepackage{booktabs}
\raggedbottom
\usepackage[hidelinks]{hyperref}

\numberwithin{equation}{section}
\newtheorem{theorem}{Theorem}[section]
\newtheorem{proposition}[theorem]{Proposition}
\newtheorem{lemma}[theorem]{Lemma}
\newtheorem{corollary}[theorem]{Corollary}
\theoremstyle{definition}
\newtheorem{definition}[theorem]{Definition}
\theoremstyle{remark}
\newtheorem{remark}[theorem]{Remark}

\newcommand{\Z}{\mathbb{Z}}
\newcommand{\Q}{\mathbb{Q}}
\newcommand{\C}{\mathbb{C}}
\newcommand{\Zl}{\Z_\ell}
\newcommand{\La}{\Lambda}
\newcommand{\Jac}{\operatorname{Jac}}
\newcommand{\coker}{\operatorname{coker}}
\newcommand{\Tor}{\operatorname{Tor}}
\newcommand{\ord}{\operatorname{ord}}
\newcommand{\car}{\operatorname{char}_{\La}}
\newcommand{\Fitt}{\operatorname{Fitt}}

\title[Algorithmic non-commutative class field theory, IV]
{Algorithmic non-commutative class field theory, IV:\\
the limit module of abelian towers\\ and the characteristic identity}

\author{Ruqing Chen}
\address{GUT Geoservice Inc., Rivi\`ere-Beaudette, Qu\'ebec, Canada}
\email{ruqing@hotmail.com}
\date{12 September 2026}
\subjclass[2020]{Primary 11R23, 05C50; Secondary 05C25, 13C12, 46L80}
\keywords{Iwasawa module, voltage graph, sandpile group, characteristic ideal,
control theorem, Bruhat--Tits tree}

\thanks{Paper IV of the series; Papers I--III are cited as [I], [II], [III]. This
paper upgrades [II, Conj.~3.6(ii-a)] from a verified growth law to a limit-module
theorem. All computational claims are reproduced by the script
\texttt{limit\_module.py} in the ancillary bundle.}

\begin{document}

\begin{abstract}
For an abelian $\Z/\ell^k$ voltage tower $(Y_k)$ over a finite base graph, Papers
II--III established the exact Iwasawa growth law of the sandpile orders and its
Newton-polygon mechanism, leaving the limit-module statement of
[II, Conj.~3.6(ii-a)] open. This paper closes it. Writing $D(t)$ for the voltage
Laplacian pencil and $\La=\Zl[[T]]$, we prove:
\[
X_\infty:=\varprojlim_k\bigl(\coker_\Z(\Delta_{Y_k})\otimes\Zl,\ \pi_*\bigr)
\;\cong\;\La^n\big/D(1+T)\,\La^n,
\]
a finitely generated torsion $\La$-module with
$\car(X_\infty)=\bigl(\det D(1+T)\bigr)$, together with an \emph{exact} control
theorem $X_\infty/\omega_kX_\infty\cong\coker_\Z(\Delta_{Y_k})\otimes\Zl$ --- exact,
not merely pseudo-isomorphic, because the tower's presentations are literally the
$\omega_k$-specializations of one $\La$-presentation. The proof rests on three
structural identities --- the derived Laplacian is the block-circulant $D(\sigma)$;
the covering pushforward is the group-ring coefficient projection; pushforward is
surjective on $\ell$-primary torsion with kernel of the exact predicted order --- each
proved and additionally machine-verified at every level of two towers. Combined with
[II, Rem.~5.4], the measured invariants $(\mu,\lambda)=(\mu(g),\lambda(g)-1)$,
$g=\det D(1+T)$, are now consequences of a theorem rather than observations. The
result is stated with explicit deference to the graph-Iwasawa literature of
Vallières, McGown--Vallières and Gonet, whose verified titles indicate substantial
overlap; what is new here, at minimum, is the exactness of the control isomorphism at
the level of Smith normal forms and the integration into the trilogy's $K$-theoretic
framework. The non-normal congruence tower ([II, Conj.~3.6(ii-b)]) and the canonicity
problem ([III, Rem.~1.10]) are untouched and remain the open frontier.
\end{abstract}

\maketitle

\section*{Status of the results}

Every numbered statement in this paper is proved in full. Lemmas~\ref{lem:block},
\ref{lem:norm} and \ref{lem:torsurj} are elementary and were additionally verified by
exact machine computation at every level $k\le3$ of the tower with voltages $(1,1,1)$,
$\ell=2$, and every level $k\le2$ of the tower $(1,0,0)$, $\ell=3$
(Table~\ref{tab:battery}). The single external algebraic input is the structure
theory of finitely generated $\La$-modules, quoted with a flagged citation. The
relation of Theorem~\ref{thm:limit} to the published graph-Iwasawa literature is
discussed, deliberately and prominently, in Section~\ref{sec:lit}.

\section{Setup and the three structural identities}\label{sec:setup}

Fix a finite connected base graph $Y_0$ on $n$ vertices, a prime $\ell$, a spanning
tree, and a voltage assignment $\alpha$ on the chords, giving the derived
$\Z/\ell^k$-covers $Y_k$ on $n\ell^k$ vertices with deck generator $\sigma$ and
covering maps $\pi\colon Y_{k+1}\to Y_k$. Let
\[
D(t)\;=\;(q+1)I_n-A(t)\ \in\ M_n\bigl(\Z[t,t^{-1}]\bigr)
\]
be the voltage Laplacian pencil ($A(t)$ the voltage adjacency), $\Delta_k$ the
ordinary Laplacian of $Y_k$, $G_k=\Z/\ell^k$, and order $V_{Y_k}$ as
$(v,g)\mapsto vN+g$, $N=\ell^k$. Write $S_N$ for the cyclic shift on $\Zl^N$ and
$Q_N\colon\Zl^{\ell N}\to\Zl^{N}$ for reduction of the group coordinate,
$e_g\mapsto e_{g\bmod N}$.

\begin{lemma}[Block-circulant identification]\label{lem:block}
In the ordering above, $\Delta_k=D(S_N)$ exactly: the level-$k$ Laplacian is the
image of the pencil under $t\mapsto S_N$. Equivalently, as a
$\Z[G_k]$-module with the deck action,
$\coker_\Z(\Delta_k)\cong\Z[G_k]^n\big/D(\sigma)\Z[G_k]^n$.
\end{lemma}

\begin{proof}
This is the definition of the derived graph read backwards: a base edge $(u,v)$ of
voltage $a$ contributes the arcs $(u,g)\!-\!(v,g{+}a)$ for all $g$, i.e.\ the block
$E_{uv}\otimes S_N^{\,a}$ plus its transpose, and the degree term is
$(q+1)I_n\otimes I_N$. Summing gives $\Delta_k=D(S_N)$; identifying
$\Zl^N\cong\Zl[G_k]$ with $S_N=$ multiplication by $\sigma$ gives the module form.
Verified as an exact matrix identity at all computed levels.
\end{proof}

\begin{lemma}[Pushforward is coefficient projection]\label{lem:norm}
Let $P_k:=I_n\otimes Q_{\ell^k}$, the matrix of the covering pushforward
$\pi_*\colon\Z^{V_{Y_{k+1}}}\to\Z^{V_{Y_k}}$ (sum over fibres). Then
\[
P_k\,\Delta_{k+1}=\Delta_k\,P_k,
\qquad
P_k\,R_{k+1}=R_k\,P_k,
\]
where $R_k=I_n\otimes S_{\ell^k}$ is the deck action; under
Lemma~\ref{lem:block} the induced map on cokernels is the group-ring coefficient
projection $\Z[G_{k+1}]\to\Z[G_k]$. Moreover $P_k$ is surjective.
\end{lemma}

\begin{proof}
Both identities reduce to the one-line facts $Q\,S_{\ell N}=S_N\,Q$ and
$Q\,S_{\ell N}^{\,a}=S_N^{\,a}\,Q$ applied blockwise to $\Delta=D(S)$; surjectivity
is immediate since every $e_{(v,h)}$ is hit. Verified as exact matrix identities at
all computed levels.
\end{proof}

\begin{lemma}[Norm surjectivity on torsion]\label{lem:torsurj}
The induced norm
$\pi_*\colon\coker_\Z(\Delta_{k+1})\to\coker_\Z(\Delta_k)$ is
surjective; it preserves the degree functional $\deg$ (total mass), hence restricts
to a surjection of torsion subgroups, with kernel of order
\[
\bigl|\ker(\pi_*|_{\Tor_\ell})\bigr|
=\ell^{\,\ord_\ell|\Jac(Y_{k+1})|\,-\,\ord_\ell|\Jac(Y_k)|}.
\]
\end{lemma}

\begin{proof}
Surjectivity of $P_k$ plus the intertwining gives surjectivity on cokernels. The
column sums of $\Delta_k$ vanish, so $\deg([y])=\sum_iy_i$ is well defined and
realizes $\coker_\Z(\Delta_k)\cong\Z\oplus\Tor$ with $\Tor=\ker(\deg)$ (a finitely
generated group surjecting onto $\Z$ with rank-zero kernel has kernel equal to its
torsion). Pushforward preserves total mass, so $\deg\circ\pi_*=\deg$; hence for
torsion $t$ downstairs any preimage $a$ has $\deg(a)=0$, i.e.\ $a$ is torsion:
surjectivity on $\Tor$. The kernel order is then forced by counting. Verified,
including the exact kernel orders $\ell^{3},\ell^{3},\ell^{3}$ and
$\ell^{1},\ell^{1}$ at the computed transitions (Table~\ref{tab:battery}).
\end{proof}

\section{The limit module and the characteristic identity}\label{sec:limit}

\begin{theorem}[Limit module]\label{thm:limit}
Let $g(T):=\det D(1+T)\in\Zl[T]$ and assume $g\neq0$ $($true in every tower of
{\rm[II, \S5]}$)$. Then, with transition maps the norms $\pi_*$ of
Lemma~\ref{lem:torsurj},
\[
X_\infty:=\varprojlim_k\bigl(\coker_\Z(\Delta_{Y_k})\otimes\Zl\bigr)
\;\cong\;\La^n\big/\,D(1+T)\,\La^n ,
\qquad \La=\Zl[[T]],
\]
under the standard identification $\varprojlim_k\Zl[G_k]\cong\La$,
$\sigma\mapsto1+T$ $($the Laurent entries $t^{\pm a}$ of the pencil landing in
$\La^\times$-multiples since $1+T\in\La^\times)$. The module $X_\infty$ is finitely
generated and torsion over $\La$, and
\[
\car(X_\infty)\;=\;\bigl(g(T)\bigr)\;=\;\bigl(\det D(1+T)\bigr).
\]
\end{theorem}

\begin{proof}
By Lemmas~\ref{lem:block}--\ref{lem:norm} the tower of cokernels is, level by level
and compatibly with the transition maps, the tower of presentations
$\Zl[G_k]^n\xrightarrow{D(\sigma)}\Zl[G_k]^n\to C_k\to0$ with transitions the
coefficient projections. All modules are finitely generated $\Zl$-modules, hence
compact, and the transition maps are continuous surjections on the presenting
copies; inverse limits are exact on such systems (Mittag--Leffler for compact
groups), so $\varprojlim C_k=\coker\bigl(\varprojlim D(\sigma)\bigr)
=\La^n/D(1+T)\La^n$. Since $g=\det D(1+T)\neq0$, the map $D(1+T)$ is injective over
the fraction field, so the cokernel is $\La$-torsion; it is finitely generated by
presentation. For the characteristic ideal, localize at a height-one prime
$P\subset\La$: $\La_P$ is a discrete valuation ring, elementary divisors give
$\operatorname{length}_{\La_P}\bigl((X_\infty)_P\bigr)=v_P\bigl(\det D(1+T)\bigr)$,
and taking the product over height-one primes yields
$\car(X_\infty)=(\det D(1+T))$ by the structure theory of finitely generated
torsion $\La$-modules \cite{Washington}.
\end{proof}

\begin{corollary}[Exact control]\label{cor:control}
For every $k\ge0$, with $\omega_k=(1+T)^{\ell^k}-1$,
\[
X_\infty\big/\omega_k X_\infty\;\cong\;\coker_\Z(\Delta_{Y_k})\otimes\Zl
\;\cong\;\bigl(\Z\oplus\Jac(Y_k)\bigr)\otimes\Zl ,
\]
an isomorphism on the nose --- not merely up to bounded kernels and cokernels ---
because $\La/\omega_k\cong\Zl[G_k]$ carries the level-$k$ presentation exactly.
\end{corollary}

\begin{corollary}[Invariants, now as a theorem]\label{cor:invariants}
The Iwasawa invariants of the tower are $\mu=\mu(g)$ and $\lambda=\lambda(g)-1$,
where $(\mu(g),\lambda(g))$ are read off the Newton polygon of $g$, the $-1$ being
the trivial-zero and normalization accounting of {\rm[II, Rem.~5.4]}. In particular
the four exact laws of {\rm[II, Table~2]} are consequences of
Theorem~\ref{thm:limit}, no longer observations.
\end{corollary}

\section{The verification battery}\label{sec:battery}

\begin{table}[ht]
\centering\small
\renewcommand{\arraystretch}{1.2}
\begin{tabular}{lccccc}
\toprule
tower & levels & (B) block & (S) deck & (P)+(C) norm identities & (N) torsion norm\\
\midrule
$(1,1,1)$, $\ell=2$ & $k\le3$ & exact & exact & exact & surj.; $|\ker|=2^3,2^3,2^3$\\
$(1,0,0)$, $\ell=3$ & $k\le2$ & exact & exact & exact & surj.; $|\ker|=3^1,3^1$\\
\bottomrule
\end{tabular}
\medskip
\caption{Machine verification of Lemmas~\ref{lem:block}--\ref{lem:torsurj}:
(B) $\Delta_k=D(S_N)$; (S) $R_k\Delta_k=\Delta_kR_k$; (P) $P_k\Delta_{k+1}
=\Delta_kP_k$; (C) $P_kR_{k+1}=R_kP_k$; (N) surjectivity of $\pi_*$ on
$\ell$-primary torsion with kernel orders exactly
$\ell^{\Delta\ord}$. Every check passed as an exact identity; the lemmas are proved
above independently, so the table is confirmation, not evidence.}
\label{tab:battery}
\end{table}

\section{Relation to the literature, stated prominently}\label{sec:lit}

The verified bibliographic records of Vallières \cite{Vallieres}, McGown--Vallières
\cite{McGownVallieres} and, above all, Gonet --- whose paper is titled precisely
\emph{Iwasawa theory of Jacobians of graphs} \cite{Gonet} --- make it very likely
that Theorem~\ref{thm:limit}, or a statement equivalent to it, is present in that
literature; the growth law certainly is, as acknowledged since [II, \S5]. Under the
series' policy, we have verified those papers' metadata but have not read their
theorems, and we therefore claim priority for nothing in
Sections~\ref{sec:setup}--\ref{sec:limit} beyond what a comparison with their texts
will leave standing. What we believe is ours at minimum: the \emph{exactness} of
Corollary~\ref{cor:control} in the form stated, its verification at the resolution
of Smith normal forms across [II, Table~2] and Table~\ref{tab:battery}, and the
integration of the statement into the trilogy's $K$-theoretic framework, where
$\coker_\Z(\Delta_{Y_k})$ is the $K_0$-group of the level-$k$ Laplacian channel
$($[I, Thm.~3.5]$)$, so that Theorem~\ref{thm:limit} reads: \emph{the
$K$-theoretic torsion of the Laplacian channel along an abelian tower is one
finitely generated torsion Iwasawa module, with characteristic ideal the voltage
pencil determinant}.

\section{What remains open}\label{sec:open}

Nothing here touches the two hard problems. The congruence tower is non-normal, its
pseudo-Iwasawa law has $\lambda=-1$, and the correct base object --- conjecturally a
module theory over the Iwahori--Hecke algebra with an intrinsic Euler-characteristic
correction ([II, Conj.~3.6(ii-b)], [III, \S2]) --- still lacks a definition of
characteristic ideal. And the canonicity problem of [III, Rem.~1.10] --- a Kirchberg
algebra over $C(\Omega)$, functorial in the Hecke data, with
$K_0\cong(\Z\oplus\Jac(Y_k))_k$ --- is exactly as open as before;
Theorem~\ref{thm:limit} sharpens its target, since the prescribed $K$-theory is now
known to be a single $\La$-module, but supplies no construction.

\section*{Acknowledgement of provenance and caveats}

Unrefereed preprint. All numbered statements are proved; the machine battery of
Table~\ref{tab:battery} confirms the three lemmas at seven levels across two towers.
External inputs: the structure theory of finitely generated torsion
$\La$-modules and the identification $\varprojlim\Zl[G_k]\cong\La$, both standard,
cited to \cite{Washington}, whose bibliographic record is now primary-source
verified (only the chapter-level pointer remains a recollection-based reading aid);
the graph-Iwasawa literature
\cite{Vallieres,McGownVallieres,Gonet}, bibliographic data verified in [II--III],
texts unread, priority explicitly ceded in Section~\ref{sec:lit}.

\begin{thebibliography}{9}
\bibitem{Vallieres} D.~Vallières, \emph{On abelian $\ell$-towers of multigraphs},
Ann.\ Math.\ Qu\'ebec \textbf{45} (2021), no.~2, 433--452. [Crossref-verified:
DOI 10.1007/s40316-020-00152-4.]
\bibitem{McGownVallieres} K.~J.~McGown and D.~Vallières, \emph{On abelian
$\ell$-towers of multigraphs II}, Ann.\ Math.\ Qu\'ebec \textbf{47}, no.~2,
461--473. [Crossref-verified: DOI 10.1007/s40316-021-00183-5; online 2021, issue
2023.]
\bibitem{Gonet} S.~R.~Gonet, \emph{Iwasawa theory of Jacobians of graphs}, Algebraic
Combinatorics \textbf{5} (2022), no.~5, 827--848. [Crossref-verified:
DOI 10.5802/alco.225.]
\bibitem{Washington} L.~C.~Washington, \emph{Introduction to Cyclotomic Fields},
2nd ed., Grad.\ Texts in Math.\ \textbf{83}, Springer, 1997; Chapter~13 for the
structure theory of $\La$-modules. [Bibliographic data primary-source verified via
the Springer book page: author, title, second edition, GTM~\textbf{83},
\copyright~1997. The chapter-level pointer is a reading aid from recollection.]
\end{thebibliography}

\end{document}
